% Imported from topics/video-generation/seedance-2.5/data-training-spec.md
\chapter{联合数据生产、噪声课程与参数组}\label{seedance-m25}
\moduleinfo{M25}{候选设计与工程选型；组合效果待验证}{\hyperref[src-ltx-2]{\texttt{ltx-2}}、\hyperref[src-ltx-original-code]{\texttt{ltx-original-code}}、\hyperref[src-tarsier2]{\texttt{tarsier2}}、\hyperref[src-tarsier-code]{\texttt{tarsier-code}}、\hyperref[src-alchemint]{\texttt{alchemint}}、\hyperref[src-vace]{\texttt{vace}}、\hyperref[src-videocof]{\texttt{videocof}}、\hyperref[src-longlive]{\texttt{longlive}}}

样本对象确定之后，继续写任务抽样、有效损失、参数组与恢复状态，明确一次更新到底包含什么。

条件已经分路以后，训练还要把它们重新接到同一损失和参数更新中。参考、编辑、几何与长叙事各有可行方法，并不表示样本一混合就能同时学好；它们的条件状态、监督有效位和梯度路径不同，也可能相互干扰，例如身份更稳但动作更静。

下面把这套联合主干的数据与更新过程写到可以检查的程度。任务比例、阈值和学习率是候选初值，用于开始实施与比较，不是已经测出的最优配方。

\section{数据生产分为八个可重做步骤}\label{seedance-m25-ux6570ux636eux751fux4ea7ux5206ux4e3aux516bux4e2aux53efux91cdux505aux6b65ux9aa4}

\begin{enumerate}
\def\labelenumi{\arabic{enumi}.}
\tightlist
\item
  媒体登记：内容哈希、来源组、可用范围、版本与音视频流信息。
\item
  时间规范化：真实重采样、PTS映射、声画偏移、色彩与旋转处理。
\item
  内容组织：镜头、场景、事件、人物与speaker候选。
\item
  关系确认：同主体跨片段配对、道具归属、speaker映射与资产家族。
\item
  来源划分：按共享主体/原片/三维资产关系确定开发、验证、测试分组。
\item
  任务派生：参考、动作、编辑、音频、几何和长叙事监督。
\item
  质量审阅：抽检与错误归因，剔除或回流重标。
\item
  编码打包：固定codec和文本特征，按形状/任务分桶，保存变换与编码版本。
\end{enumerate}

来源分组还会随新数据改变。若新参考把原来两个来源组连起来，就要重新检查 split；测试身份已经固定时，也不能把它的新视角加进训练，再继续声称身份隔离。去重关系因此需要随数据版本一起维护。

\section{每种自动工具输出什么}\label{seedance-m25-ux6bcfux79cdux81eaux52a8ux5de5ux5177ux8f93ux51faux4ec0ux4e48}

{\def\LTcaptype{table} % do not increment counter
\small\setlength{\tabcolsep}{3pt}\renewcommand{\arraystretch}{1.18}
\begin{longtable}[]{@{}
  >{\raggedright\arraybackslash}p{(\linewidth - 4\tabcolsep) * \real{0.1900}}
  >{\raggedright\arraybackslash}p{(\linewidth - 4\tabcolsep) * \real{0.3900}}
  >{\raggedright\arraybackslash}p{(\linewidth - 4\tabcolsep) * \real{0.4200}}@{}}
\toprule\noalign{}
\begin{minipage}[b]{\linewidth}\raggedright
环节
\end{minipage} & \begin{minipage}[b]{\linewidth}\raggedright
可选工具/来源
\end{minipage} & \begin{minipage}[b]{\linewidth}\raggedright
必须保存的结果
\end{minipage} \\
\midrule\noalign{}
\endhead
\bottomrule\noalign{}
\endlastfoot
媒体探测与转码 & ffprobe、FFmpeg & 原时间戳、变换参数、输出映射 \\
镜头与场景组织 & 场景检测器＋连续性审阅 & 切点、误切标记、场景组 \\
视频描述 & Tarsier/Tarsier2等 & 事实描述、事件时间、模型版本与人工纠错 \\
主体检测与跟踪 & Grounding DINO、SAM2等 & 框、逐帧mask、轨迹、可见性与置信度 \\
语音转录与分段 & ASR、diarization与必要的声源分离 & 台词、speaker段、重叠讲话和声画偏移 \\
结构条件 & 3D渲染器或深度/姿态/相机估计器 & 单位、坐标、有效位、估计置信度 \\
\end{longtable}
}

工具是候选生产组件，名称不代表已经安装或运行。阈值先在开发集按人工错误类型校准并登记版本；评估集上不再反复调阈值。

\section{八个数据池如何构造监督}\label{seedance-m25-ux516bux4e2aux6570ux636eux6c60ux5982ux4f55ux6784ux9020ux76d1ux7763}

{\def\LTcaptype{table} % do not increment counter
\small\setlength{\tabcolsep}{3pt}\renewcommand{\arraystretch}{1.18}
\begin{longtable}[]{@{}
  >{\raggedright\arraybackslash}p{(\linewidth - 4\tabcolsep) * \real{0.1900}}
  >{\raggedright\arraybackslash}p{(\linewidth - 4\tabcolsep) * \real{0.3900}}
  >{\raggedright\arraybackslash}p{(\linewidth - 4\tabcolsep) * \real{0.4200}}@{}}
\toprule\noalign{}
\begin{minipage}[b]{\linewidth}\raggedright
池
\end{minipage} & \begin{minipage}[b]{\linewidth}\raggedright
配对方式
\end{minipage} & \begin{minipage}[b]{\linewidth}\raggedright
不应混入的伪真值
\end{minipage} \\
\midrule\noalign{}
\endhead
\bottomrule\noalign{}
\endlastfoot
基础AV & 同一原片中的caption、视频和对齐音轨 & 猜测的画外台词、错位声音 \\
身份参考 & 同主体不同时间/视角参考与目标 & 只靠相似embedding认定身份 \\
动作/相机 & 目标提取结构，或同轨迹跨外观重渲染 & 用原人物外观暗中帮助控制 \\
声音参考 & 同speaker不同台词的干净语音与目标 & 背景音乐被标为音色 \\
区域编辑 & 人工编辑对、受控重渲染、明确标记的弱监督 & 把相机变化误标为单一颜色变化 \\
示例传播 & 源关键帧、编辑后关键帧与可靠目标视频 & 每帧独立图像编辑产生的闪烁``真值'' \\
几何控制 & 同资产与相机的控制层和目标材质渲染 & 米制渲染和未知尺度估计混用 \\
场景叙事 & 连续片段、镜头事件和角色道具状态 & 无关短片拼接成虚假连续场景 \\
\end{longtable}
}

难例也要主动收集：参考与目标相隔较远、同人换衣、中途缺席、动作换身份，以及 speaker 重叠。它们分别检验条件的不同含义。故意错绑参考的样本改用于冲突检测或比较，不继续拿原目标作普通 flow 真值。

先为每池制作一个分层审阅批次，例如200条、其中50条双人独立标注。这是流程校准数量。记录配对、时间、mask、caption和speaker各自错误率；没有足够真实样本时，减少解锁能力，而不是凭空补足池大小。

\section{样本有三个独立状态轴}\label{seedance-m25-ux6837ux672cux6709ux4e09ux4e2aux72ecux7acbux72b6ux6001ux8f74}

每个模态分别记录是否存在、是否算隐藏状态、是否更新 latent，每个目标再记录监督有效位。全局 batch 尽量按相同任务、形状和声音状态组织，是为了减少 padding 与分布式未使用参数的歧义，不是说所有样本天生使用同一损失。

\begin{verbatim}
基础AV：两流均存在、处理、更新，计算两流loss。
无音轨视频：只有视频目标；音频不存在，跳过声音与无效交互。
保留声音的编辑：音频存在且处理，但不更新latent、不计算音频loss。
局部对白编辑：活动声音区加噪并监督，外部为干净条件。
因果续写：历史为干净条件，只监督当前块，两流同一块划分。
\end{verbatim}

缺音轨没有声音目标，静音则有；保留声音又是另一种状态，它固定音频输出，却仍可能需要声音隐藏流提供口型条件。先区分这些变量，才能知道哪一项损失应当存在。

\section{噪声课程不能盲目按长序列外推}\label{seedance-m25-ux566aux58f0ux8bfeux7a0bux4e0dux80fdux76f2ux76eeux6309ux957fux5e8fux5217ux5916ux63a8}

原系列训练sampler使用长度相关的shifted logit-normal。它的偏移为：

\[
\mu(N)=0.95+\frac{1.10}{3072}(N-1024),
\qquad \sigma=\operatorname{sigmoid}(\xi+\mu(N)),\quad
\xi\sim\mathcal N(0,1).
\]

原实现按实际长度外推，并没有在4096处自动截断。当 \(N=40\,768\) 时，\(\mu\approx15.18\)，分布中位数接近0.9999997。长阶段照抄会高度集中在几乎纯噪声的位置，缺少足够的中低噪声训练。

候选长阶段采用：90\%有界shifted logit-normal、10\%均匀sigma；偏移限制在0.95至2.05之间。长度只计算目标视频token，不加参考摘要数；短片原生配置单独保留为对照。推理sigma调度也固定同一候选shift策略与原生端点处理，避免训练改了、验证仍使用另一套外推。

这个课程明确改变了原配方，所以必须保留原生对照，不能先叫作无损改进。基础 AV 两流共用 sigma，固定条件为 0，异步镜头另抽镜头 sigma；首版暂不增加没有依据的任意双流独立组合。

\section{损失先按样本有效元素平均}\label{seedance-m25-ux635fux5931ux5148ux6309ux6837ux672cux6709ux6548ux5143ux7d20ux5e73ux5747}

对样本 \(b\)、模态 \(m\)，先在FP32中计算：

\[
\ell_b^m=\frac{\sum_i W_{bi}^m
(u_{bi}^m-u_{bi}^{m*})^2}{\sum_i W_{bi}^m},
\qquad u^{m*}=\epsilon^m-z_0^m.
\]

分子累加真正参与训练的标量误差，分母累计同一组标量的有效权重，所以求得的是每样本均值。缺少目标时分母为零，含义是该模态没有监督，而不是一个需要用极小常数掩盖的数值问题。

\(i\)遍历标量通道元素，token mask需广播到通道。没有有效目标的样本不进入该模态平均；不能给零分母加epsilon后把缺失监督当作一个低loss样本。

平均的顺序也定义了训练对象。先在每个样本内归一，再对同任务的有效样本平均，长样本就不会只因 token 多而自动占更大权重。多卡有效数量不同时，应归约全局计数并修正梯度缩放，不能再简单平均各卡均值。

基础目标为 \(\ell_v+\lambda_a\ell_a\)，初始 \(\lambda_a=1\)。区域头使用BCE＋Dice或独立区域flow，选择一种输出形式；保持项初始权重0.1且仅用于适用的对齐编辑。极简有效位归一化见 \hyperref[code-masked-sample-loss]{masked\_sample\_loss.py}。

\section{参数组与解冻顺序}\label{seedance-m25-ux53c2ux6570ux7ec4ux4e0eux89e3ux51bbux987aux5e8f}

{\def\LTcaptype{table} % do not increment counter
\small\setlength{\tabcolsep}{3pt}\renewcommand{\arraystretch}{1.18}
\begin{longtable}[]{@{}
  >{\raggedright\arraybackslash}p{(\linewidth - 4\tabcolsep) * \real{0.1900}}
  >{\raggedright\arraybackslash}p{(\linewidth - 4\tabcolsep) * \real{0.3900}}
  >{\raggedleft\arraybackslash}p{(\linewidth - 4\tabcolsep) * \real{0.4200}}@{}}
\toprule\noalign{}
\begin{minipage}[b]{\linewidth}\raggedright
阶段
\end{minipage} & \begin{minipage}[b]{\linewidth}\raggedright
更新对象
\end{minipage} & \begin{minipage}[b]{\linewidth}\raggedleft
首版候选学习率
\end{minipage} \\
\midrule\noalign{}
\endhead
\bottomrule\noalign{}
\endlastfoot
支路起步 & 池化器、绑定表示、条件塔、投影与门控 & \(10^{-4}\) \\
编辑联合 & 已有支路＋源条件与区域模块 & 新模块 \(10^{-4}\)，已有模块 \(3\times10^{-5}\) \\
主干适配 & 上述模块＋选定原生注意力LoRA & LoRA \(10^{-5}\) \\
因果派生 & 因果适配器、必要LoRA与当前块策略 & 重新短pilot确定，不继承错误的全局sigma \\
\end{longtable}
}

codec、vocoder和文本模型首版冻结。LoRA起点rank16、alpha16、A正常初始化而B置零；先只匹配原生self与AV注意力的 \codepath{to_q}、\codepath{to_v}，排除文本注意力和新增支路，再登记实际模块全名及参数数。通用后缀 \codepath{to_q} 可能匹配所有分支，不能直接作为精确参数组。

LoRA 或新增支路参数较少，但损失到它们的梯度仍可能经过很多冻结块。冻结权重只是不更新它，并未取消这段导数。把整个主干放入 \codepath{no_grad()} 会截断路径，梯度检查点减少 activation，也没有让主干反传成本消失。

候选有效batch8、全局范数裁剪1.0、100步预热、每100次更新固定验证；先1000次更新做短pilot。AdamW的weight decay对bias、norm和门控单列为0，普通投影初始0.01。所有数值都是实施起点。

\section{任务混合与干扰处理}\label{seedance-m25-ux4efbux52a1ux6df7ux5408ux4e0eux5e72ux6270ux5904ux7406}

支路阶段按20\%基础AV、30\%图片参考、20\%动作/相机、20\%声音参考、10\%几何采样；编辑阶段按15\%基础、20\%参考、30\%编辑对、20\%示例传播、15\%区域定位采样。缺某池时重新发布混合版本，不悄悄沿用一个无法实现的比例。

每次全局更新先按比例抽任务，再抽来源组与形状桶。参考数从1--2到4/8，再到容量上限；时长从8到16/30秒。不要同时改变时长、分辨率、参考数和可训练模块，导致无法定位退化。

出现退化后先按任务看独立 loss、固定切片与共享参数梯度。身份提高但动作趋静，检查是否复制数据或条件过强；音频压过视频，检查有效元素与梯度比例；编辑变成复制，再看目标变化，不只看保持项下降。这样观察才对应得上前面定义的未知量。

\section{checkpoint是一个组合对象}\label{seedance-m25-checkpointux662fux4e00ux4e2aux7ec4ux5408ux5bf9ux8c61}

一个恢复点要能说明这一步怎样得到：基座、新模块、LoRA、参数组、optimizer、scheduler、随机状态、采样游标、数据、任务混合、sigma 与条件预算，都要一起保存。只存 LoRA 可能丢条件塔，只存模型又丢掉更新状态，均不能证明恢复了同一训练轨迹。

因此选定 trainer 后还要补训练状态的保存与恢复。原函数主要写模型权重，LoRA 路径又经 PEFT 筛选，未自动满足上面的要求。接入新模块时，反向检查优化器参数组和保存清单是否对应，避免训练确实更新了，交付却没有这份权重。

真实训练接口读 \href{https://github.com/Lightricks/LTX-2/blob/bd92a5f408eb8c99f8c05e9fe2fcaaec8b56b638/packages/ltx-trainer/src/ltx_trainer/training_strategies/text_to_video.py}{基础AV策略}、\href{https://github.com/Lightricks/LTX-2/blob/bd92a5f408eb8c99f8c05e9fe2fcaaec8b56b638/packages/ltx-trainer/src/ltx_trainer/timestep_samplers.py}{timestep samplers} 与 \href{https://github.com/Lightricks/LTX-2/blob/bd92a5f408eb8c99f8c05e9fe2fcaaec8b56b638/packages/ltx-trainer/src/ltx_trainer/trainer.py}{trainer}。本协议规定新增工作如何组织，没有把纸面候选值当成已收敛训练结果。
